\documentclass[11pt,a4paper]{ctexart}

\usepackage[margin=2.2cm]{geometry}
\usepackage{amsmath,amssymb,amsthm}
\usepackage{enumitem}
\usepackage{mathtools}

\setlength{\parindent}{0pt}
\setlength{\parskip}{0.5em}
\setlist{leftmargin=2em,itemsep=0.3em}

\newenvironment{problem}[1]
  {\par\medskip\noindent\textbf{#1}\quad}
  {\par}
\newenvironment{solution}
  {\par\noindent\textbf{解}\quad}
  {\par}
\renewcommand{\proofname}{证明}

\title{线性代数（MATH1205）第二次作业}
\author{}
\date{}

\begin{document}
\maketitle
\vspace{-3em}

% ============================================================
\section*{第 1 题}

\begin{problem}{}设
$A=\begin{pmatrix}1&3\\-1&0\\2&3\end{pmatrix}$，
$B=\begin{pmatrix}-2&0\\2&-1\\-1&1\end{pmatrix}$，
$C=\begin{pmatrix}-1&2\\2&1\\2&-3\end{pmatrix}$，
求 $A+B$，$A-B+C$，$-A+3B+2C$。
\end{problem}

\begin{solution}
矩阵的加法与数乘都是对应元素逐个运算。
\[
A+B=\begin{pmatrix}1-2&3+0\\-1+2&0-1\\2-1&3+1\end{pmatrix}
=\begin{pmatrix}-1&3\\1&-1\\1&4\end{pmatrix}.
\]
\[
A-B+C=\begin{pmatrix}1+2-1&3-0+2\\-1-2+2&0+1+1\\2+1+2&3-1-3\end{pmatrix}
=\begin{pmatrix}2&5\\-1&2\\5&-1\end{pmatrix}.
\]
\[
-A+3B+2C
=\begin{pmatrix}-1&-3\\1&0\\-2&-3\end{pmatrix}
+\begin{pmatrix}-6&0\\6&-3\\-3&3\end{pmatrix}
+\begin{pmatrix}-2&4\\4&2\\4&-6\end{pmatrix}
=\begin{pmatrix}-9&1\\11&-1\\-1&-6\end{pmatrix}.
\]
\end{solution}

% ============================================================
\section*{第 2 题（原题号 3）}

\begin{problem}{}设
$A=\begin{pmatrix}x&0\\7&y\end{pmatrix}$，
$B=\begin{pmatrix}u&2v\\y&2\end{pmatrix}$，
$C=\begin{pmatrix}3&-4\\x&v\end{pmatrix}$，
且 $A+2B=C$，求 $x,y,u,v$。
\end{problem}

\begin{solution}
\[
A+2B=\begin{pmatrix}x+2u&4v\\7+2y&y+4\end{pmatrix}
=\begin{pmatrix}3&-4\\x&v\end{pmatrix}.
\]
比较对应元素得
\[
\begin{cases}
x+2u=3,\\ 4v=-4,\\ 7+2y=x,\\ y+4=v.
\end{cases}
\]
由第二式 $v=-1$；代入第四式 $y=v-4=-5$；代入第三式 $x=7+2\cdot(-5)=-3$；
代入第一式 $2u=3-x=6$，$u=3$。

所以 $x=-3,\ y=-5,\ u=3,\ v=-1$。
\end{solution}

% ============================================================
\section*{第 3 题（原题号 4）}

\begin{problem}{}设 $A=(a_{ij})_{m\times n}$，$k$ 为数。证明：若 $kA=O$，则 $k=0$ 或 $A=O$。\end{problem}

\begin{proof}
由数乘定义，$kA=(ka_{ij})_{m\times n}$。$kA=O$ 意味着
\[
ka_{ij}=0,\qquad i=1,\dots,m,\ j=1,\dots,n.
\]
若 $k=0$，结论已成立。若 $k\neq0$，则 $k$ 可逆（$k^{-1}$ 存在），在每个等式两边乘 $k^{-1}$ 得
\[
a_{ij}=k^{-1}(ka_{ij})=k^{-1}\cdot0=0,\qquad \text{对所有 } i,j,
\]
即 $A=O$。故 $k=0$ 或 $A=O$。
\end{proof}

% ============================================================
\section*{第 4 题（原题号 5）}

\begin{problem}{}计算下列矩阵乘积。\end{problem}

\textbf{(2)}
\[
\begin{pmatrix}1&3&-1\\0&4&2\\7&0&1\end{pmatrix}
\begin{pmatrix}1\\-2\\3\end{pmatrix}
=\begin{pmatrix}1\cdot1+3\cdot(-2)+(-1)\cdot3\\0\cdot1+4\cdot(-2)+2\cdot3\\7\cdot1+0\cdot(-2)+1\cdot3\end{pmatrix}
=\begin{pmatrix}-8\\-2\\10\end{pmatrix}.
\]

\textbf{(4)} $3\times1$ 矩阵乘 $1\times3$ 矩阵，结果为 $3\times3$ 矩阵，第 $i$ 行第 $j$ 列为 $a_i b_j$：
\[
\begin{pmatrix}1\\-2\\3\end{pmatrix}(3,-1,2)
=\begin{pmatrix}3&-1&2\\-6&2&-4\\9&-3&6\end{pmatrix}.
\]

% ============================================================
\section*{第 5 题（原题号 6）}

\begin{problem}{}证明两个同阶上三角矩阵的乘积还是上三角矩阵。\end{problem}

\begin{proof}
设 $A=(a_{ij})$，$B=(b_{ij})$ 都是 $n$ 阶上三角矩阵，即当 $i>j$ 时 $a_{ij}=0$，$b_{ij}=0$。
记 $C=AB=(c_{ij})$，则
\[
c_{ij}=\sum_{k=1}^{n}a_{ik}b_{kj}.
\]
取 $i>j$，考察和式中的每一项 $a_{ik}b_{kj}$：
\begin{itemize}
  \item 若 $k<i$，则 $a_{ik}=0$；
  \item 若 $k\ge i$，则 $k\ge i>j$，从而 $b_{kj}=0$。
\end{itemize}
所以每一项都是 $0$，$c_{ij}=0$（$i>j$）。因此 $C=AB$ 是上三角矩阵。
\end{proof}

% ============================================================
\section*{第 6 题（原题号 8）}

\begin{problem}{}计算下列各题。\end{problem}

\textbf{(1)} $\begin{pmatrix}1&0\\\lambda&1\end{pmatrix}^n$。

\begin{solution}
断言 $\begin{pmatrix}1&0\\\lambda&1\end{pmatrix}^n=\begin{pmatrix}1&0\\n\lambda&1\end{pmatrix}$。
$n=1$ 时显然。设 $n$ 时成立，则
\[
\begin{pmatrix}1&0\\\lambda&1\end{pmatrix}^{n+1}
=\begin{pmatrix}1&0\\n\lambda&1\end{pmatrix}\begin{pmatrix}1&0\\\lambda&1\end{pmatrix}
=\begin{pmatrix}1&0\\n\lambda+\lambda&1\end{pmatrix}
=\begin{pmatrix}1&0\\(n+1)\lambda&1\end{pmatrix}.
\]
由数学归纳法，
$\begin{pmatrix}1&0\\\lambda&1\end{pmatrix}^n=\begin{pmatrix}1&0\\n\lambda&1\end{pmatrix}$。
\end{solution}

\textbf{(2)} $\begin{pmatrix}\cos\theta&-\sin\theta\\\sin\theta&\cos\theta\end{pmatrix}^n$。

\begin{solution}
记 $R(\theta)=\begin{pmatrix}\cos\theta&-\sin\theta\\\sin\theta&\cos\theta\end{pmatrix}$。由和角公式，
\[
R(\alpha)R(\beta)
=\begin{pmatrix}
\cos\alpha\cos\beta-\sin\alpha\sin\beta & -(\cos\alpha\sin\beta+\sin\alpha\cos\beta)\\
\sin\alpha\cos\beta+\cos\alpha\sin\beta & \cos\alpha\cos\beta-\sin\alpha\sin\beta
\end{pmatrix}
=R(\alpha+\beta).
\]
于是对 $n$ 归纳：$n=1$ 显然；若 $R(\theta)^n=R(n\theta)$，则 $R(\theta)^{n+1}=R(n\theta)R(\theta)=R((n+1)\theta)$。所以
\[
\begin{pmatrix}\cos\theta&-\sin\theta\\\sin\theta&\cos\theta\end{pmatrix}^n
=\begin{pmatrix}\cos n\theta&-\sin n\theta\\\sin n\theta&\cos n\theta\end{pmatrix}.
\]
\end{solution}

\textbf{(3)} $(\alpha\beta^T)^n$，$\alpha=\begin{pmatrix}a_1\\a_2\\a_3\end{pmatrix}$，$\beta=\begin{pmatrix}b_1\\b_2\\b_3\end{pmatrix}$；
当 $\alpha=\begin{pmatrix}1\\2\\3\end{pmatrix}$，$\beta=\begin{pmatrix}1\\\frac12\\\frac13\end{pmatrix}$ 时写出具体结果。

\begin{solution}
$\beta^T\alpha=a_1b_1+a_2b_2+a_3b_3$ 是一个数，记为 $s$。由矩阵乘法结合律，
\[
(\alpha\beta^T)^n=\alpha(\beta^T\alpha)(\beta^T\alpha)\cdots(\beta^T\alpha)\beta^T
=s^{\,n-1}\,\alpha\beta^T,
\]
即
\[
(\alpha\beta^T)^n=(a_1b_1+a_2b_2+a_3b_3)^{n-1}
\begin{pmatrix}a_1b_1&a_1b_2&a_1b_3\\a_2b_1&a_2b_2&a_2b_3\\a_3b_1&a_3b_2&a_3b_3\end{pmatrix}.
\]
当 $\alpha=(1,2,3)^T$，$\beta=(1,\frac12,\frac13)^T$ 时，$s=1+1+1=3$，
\[
\alpha\beta^T=\begin{pmatrix}1&\frac12&\frac13\\2&1&\frac23\\3&\frac32&1\end{pmatrix},
\qquad
(\alpha\beta^T)^n=3^{n-1}\begin{pmatrix}1&\frac12&\frac13\\2&1&\frac23\\3&\frac32&1\end{pmatrix}.
\]
\end{solution}

\textbf{(4)} $\begin{pmatrix}1&-1&-1&-1\\-1&1&-1&-1\\-1&-1&1&-1\\-1&-1&-1&1\end{pmatrix}^n$。

\begin{solution}
记该矩阵为 $A$。直接计算
\[
A^2=\begin{pmatrix}4&0&0&0\\0&4&0&0\\0&0&4&0\\0&0&0&4\end{pmatrix}=4E.
\]
（例如第 $1$ 行乘第 $1$ 列：$1+1+1+1=4$；第 $1$ 行乘第 $2$ 列：$-1-1+1+1=0$；其余同理。）
于是 $A^{2k}=(A^2)^k=4^kE$，$A^{2k+1}=4^kA$，即
\[
A^n=\begin{cases}
4^{n/2}E=2^nE, & n \text{ 为偶数},\\[4pt]
4^{(n-1)/2}A=2^{n-1}A, & n \text{ 为奇数}.
\end{cases}
\]
\end{solution}

\textbf{(5)} 设 $A=\begin{pmatrix}1&0&1\\0&2&0\\1&0&1\end{pmatrix}$，求 $A^n-2A^{n-1}$（$n\ge2$）。

\begin{solution}
\[
A^2=\begin{pmatrix}1&0&1\\0&2&0\\1&0&1\end{pmatrix}\begin{pmatrix}1&0&1\\0&2&0\\1&0&1\end{pmatrix}
=\begin{pmatrix}2&0&2\\0&4&0\\2&0&2\end{pmatrix}=2A,
\]
即 $A^2-2A=O$。于是当 $n\ge2$ 时（约定 $A^0=E$），
\[
A^n-2A^{n-1}=A^{n-2}(A^2-2A)=A^{n-2}O=O.
\]
\end{solution}

% ============================================================
\section*{第 7 题（原题号 9）}

\begin{problem}{}求与矩阵 $N=\begin{pmatrix}0&1&0&0\\0&0&1&0\\0&0&0&1\\0&0&0&0\end{pmatrix}$ 可交换的所有矩阵。\end{problem}

\begin{solution}
设 $X=(x_{ij})_{4\times4}$ 满足 $NX=XN$。左乘 $N$ 把 $X$ 的各行上移一行，右乘 $N$ 把 $X$ 的各列右移一列：
\[
NX=\begin{pmatrix}
x_{21}&x_{22}&x_{23}&x_{24}\\
x_{31}&x_{32}&x_{33}&x_{34}\\
x_{41}&x_{42}&x_{43}&x_{44}\\
0&0&0&0
\end{pmatrix},\qquad
XN=\begin{pmatrix}
0&x_{11}&x_{12}&x_{13}\\
0&x_{21}&x_{22}&x_{23}\\
0&x_{31}&x_{32}&x_{33}\\
0&x_{41}&x_{42}&x_{43}
\end{pmatrix}.
\]
比较对应元素：
\begin{itemize}
  \item 第 $1$ 列：$x_{21}=x_{31}=x_{41}=0$；
  \item 第 $4$ 行：$x_{41}=x_{42}=x_{43}=0$；
  \item 其余位置：$x_{22}=x_{11}$，$x_{23}=x_{12}$，$x_{24}=x_{13}$，
        $x_{32}=x_{21}=0$，$x_{33}=x_{22}$，$x_{34}=x_{23}$，
        $x_{42}=x_{31}$，$x_{43}=x_{32}$，$x_{44}=x_{33}$。
\end{itemize}
所以主对角线以下全为 $0$，且每条平行于主对角线的斜线上元素相等：
$x_{11}=x_{22}=x_{33}=x_{44}=a$，$x_{12}=x_{23}=x_{34}=b$，$x_{13}=x_{24}=c$，$x_{14}=d$。
与 $N$ 可交换的矩阵为
\[
X=\begin{pmatrix}a&b&c&d\\0&a&b&c\\0&0&a&b\\0&0&0&a\end{pmatrix}
=aE+bN+cN^2+dN^3,
\]
其中 $a,b,c,d$ 为任意数。反之，这样的 $X$ 是 $N$ 的多项式，显然与 $N$ 可交换。
\end{solution}

% ============================================================
\section*{第 8 题（原题号 10）}

\begin{problem}{}求一个满足 $A^2=O$ 且 $A\neq O$ 的二阶方阵 $A_{2\times2}$。\end{problem}

\begin{solution}
取 $A=\begin{pmatrix}0&1\\0&0\end{pmatrix}\neq O$，则
\[
A^2=\begin{pmatrix}0&1\\0&0\end{pmatrix}\begin{pmatrix}0&1\\0&0\end{pmatrix}
=\begin{pmatrix}0\cdot0+1\cdot0&0\cdot1+1\cdot0\\0&0\end{pmatrix}=O.
\]
（一般地，$\begin{pmatrix}a&b\\c&-a\end{pmatrix}$ 满足 $a^2+bc=0$ 且不全为零时都符合要求。）
\end{solution}

% ============================================================
\section*{第 9 题（原题号 11）}

\begin{problem}{}设 $A=\begin{pmatrix}1&-1\\2&3\end{pmatrix}$，求：(1) $A^2-2A$；(2) $3A^3-2A^2+5A-4E$。\end{problem}

\begin{solution}
\textbf{(1)}
\[
A^2=\begin{pmatrix}1&-1\\2&3\end{pmatrix}\begin{pmatrix}1&-1\\2&3\end{pmatrix}
=\begin{pmatrix}1-2&-1-3\\2+6&-2+9\end{pmatrix}
=\begin{pmatrix}-1&-4\\8&7\end{pmatrix},
\]
\[
A^2-2A=\begin{pmatrix}-1&-4\\8&7\end{pmatrix}-\begin{pmatrix}2&-2\\4&6\end{pmatrix}
=\begin{pmatrix}-3&-2\\4&1\end{pmatrix}.
\]

\textbf{(2)}
\[
A^3=A^2A=\begin{pmatrix}-1&-4\\8&7\end{pmatrix}\begin{pmatrix}1&-1\\2&3\end{pmatrix}
=\begin{pmatrix}-1-8&1-12\\8+14&-8+21\end{pmatrix}
=\begin{pmatrix}-9&-11\\22&13\end{pmatrix}.
\]
于是
\[
\begin{aligned}
3A^3-2A^2+5A-4E
&=\begin{pmatrix}-27&-33\\66&39\end{pmatrix}
+\begin{pmatrix}2&8\\-16&-14\end{pmatrix}
+\begin{pmatrix}5&-5\\10&15\end{pmatrix}
+\begin{pmatrix}-4&0\\0&-4\end{pmatrix}\\
&=\begin{pmatrix}-24&-30\\60&36\end{pmatrix}.
\end{aligned}
\]
\end{solution}

\end{document}
